\subsubsection{文件目录操作}
在src文件夹中的dir.c文件中实现多进程下访问目录功能,这里显然创建目录、列出当前目录下所有文件及文件夹操作都需要互斥进行,具体代码如下:
\begin{ccode}{目录操作函数}
#include "fat.h"

bool mkdir(int pid, const char *dirname)
{
    if (pid < 0 || pid >= MAX_PROCESS) return false;
    int cur_dir = ProcessTable[pid].current_dir;

    pthread_mutex_lock(&fs_mutex);//创建文件需要加锁

    for (int i = 0; i < MAX_ENTRY; i++) {
        if (Directory[i].ac == 0) {
            strncpy(Directory[i].name, dirname, 16);
            Directory[i].size = 0;
            Directory[i].start_block = -1;
            Directory[i].ac = 1;
            Directory[i].type = 1;
            Directory[i].parent = cur_dir;
            Directory[i].reader_count = 0;
            sem_init(&Directory[i].rw_lock, 0, 1);
            pthread_mutex_init(&Directory[i].mutex, NULL);
            pthread_mutex_unlock(&fs_mutex);
            return true;
        }
    }

    pthread_mutex_unlock(&fs_mutex);
    return false;
}

static int get_dir_size(int dir_index)
{
    int total = 0;
    for (int i = 0; i < MAX_ENTRY; i++) {
        if (Directory[i].ac == 1 && Directory[i].parent == dir_index) {
            if (Directory[i].type == 0) {
                total += Directory[i].size;
            } else if (Directory[i].type == 1) {
                total += get_dir_size(i);
            }
        }
    }
    return total;
}

bool ls(int pid)
{
    if (pid < 0 || pid >= MAX_PROCESS) return false;
    int cur_dir = ProcessTable[pid].current_dir;

    pthread_mutex_lock(&fs_mutex);

    for (int i = 0; i < MAX_ENTRY; i++) {
        if (Directory[i].ac == 1 && Directory[i].parent == cur_dir) {
            int show_size;
            if (Directory[i].type == 1) {
                show_size = get_dir_size(i);
            } else {
                show_size = Directory[i].size;
            }
            printf("%s\t%s\t%d bytes\n",
                   Directory[i].type == 1 ? "DIR" : "FILE",
                   Directory[i].name,
                   show_size);
        }
    }

    pthread_mutex_unlock(&fs_mutex);
    return true;
}

bool cd(int pid, const char *dirname)
{
    if (pid < 0 || pid >= MAX_PROCESS || dirname == NULL) return false;

    int cur_dir = ProcessTable[pid].current_dir;

    if (strcmp(dirname, "/") == 0) {
        ProcessTable[pid].current_dir = 0;
        return true;
    }

    if (strcmp(dirname, "..") == 0) {
        if (Directory[cur_dir].parent != -1) {
            ProcessTable[pid].current_dir = Directory[cur_dir].parent;
        }
        return true;
    }

    //遍历以在当前目录下找子目录(对Directory[]只读),避免其他进程此时mkdir/touch,加锁保证一致性
    pthread_mutex_lock(&fs_mutex);

    for (int i = 0; i < MAX_ENTRY; i++) {
        if (Directory[i].ac == 1 &&
            Directory[i].parent == cur_dir &&
            Directory[i].type == 1 &&
            strcmp(Directory[i].name, dirname) == 0) {
            ProcessTable[pid].current_dir = i;
            pthread_mutex_unlock(&fs_mutex);
            return true;
        }
    }

    pthread_mutex_unlock(&fs_mutex);
    return false;
}
\end{ccode}
目录操作有关函数首先需要提取进程控制块中存放的当前目录位置作为实际工作目录。mkdir()函数在查找空闲目录项前需要申请信号量(pthread\_mutex\_lock(\&fs\_mutex)),查找完成后成功
获取空闲Directory项作为新建目录项或查找失败,都需要释放信号量(pthread\_mutex\_unlock(\&fs\_mutex)),其中成功获取空闲项并修改的操作与任务一相同。类似地,ls()函数和cd()函数的核心操作与任务一中类似,
只是在需要互斥访问目录项数组Directory时需要前后加信号量保护。